% ---
% id: receiver/rem:acyc-efficiency
% kind: rem
% title: Where the element graph enters the efficiency theorem
% summary: >-
%   The claims that read only the Gram substitution of the efficiency theorem
%   involve the terminal efficiency layer alone, a function of the port map and
%   efficiencies, and no element graph. The theorem's items on the compiled
%   product and its record hold only for a well-formed, hence acyclic, graph.
% prominence: standard
% status: accepted
% uses: []
% mentions: [def:compile, def:graph, receiver/detector-imperfections/thm:efficiency, receiver/detector-imperfections/rem:powervsclick, receiver/detector-imperfections/rem:sperlingeta, receiver/detector-imperfections/req:efficiency-convention, receiver/detector-imperfections/req:efficiency-schema, receiver/detector-imperfections/req:efficiency-record]
% assumes: []
% refs: []
% provenance: newly-derived
% lean:
% ---

\begin{remark}
\label{rem:acyc-efficiency}
The two remarks that use Theorem~\ref{thm:efficiency}
(Remarks~\ref{rem:powervsclick} and~\ref{rem:sperlingeta}) and
Requirement~\ref{req:efficiency-convention}, which rests on the first of them,
read only the substitution of item (i) of that theorem: they involve the
\tref{terminal detector-efficiency layer} $D_\eta$ of
Definition~\ref{def:compile}, which is a function of the port map and the
per-port efficiencies alone, and no graph enters their statements or proofs.
The same holds of items (i), (ii) and (v) of the theorem. Items (iii) and (iv)
are statements about the compiled product of Definition~\ref{def:compile} and
its record, which exist only for a well-formed, hence acyclic, graph
(condition (E2) of Definition~\ref{def:graph}); so are
Requirements~\ref{req:efficiency-schema} and~\ref{req:efficiency-record},
which state their engine side.
\end{remark}
